\section{后记：之后的数学}
\label{sec:ch10-epilogue}
% ============================================================================

写到这里，
全书最戏剧性的故事已经结束：
一条假想的 Frey 曲线被模性与水平下降夹击，
最终在空空间 $\Sk_2(\GammaO(2))$ 前消失。
但数学生命力最强的地方，
常常恰恰在“定理证完以后”。
Wiles 的工作并没有把这个领域封口，
相反，它把一整片新大陆真正打开了。

\textbf{我们遇到了什么问题？}
费马大定理虽然已经解决，
可它背后的方法——模性、Galois 表示、形变理论、局部--整体对应——
远不止服务于一个古老方程。
人们自然会追问：
既然半稳定椭圆曲线都模，
那么任意椭圆曲线呢？
更一般的 Abel 簇呢？
更高维的 Galois 表示呢？
所有这些问题最终都汇聚到更大的背景：Langlands 纲领。

\textbf{核心思想是什么？}
谷山--志村定理之所以如此重要，
不是因为它只解决了 FLT，
而是因为它展示了一个范例：
一个几何对象的 $L$-函数，
可以由一个自守对象来描述；
一个 Galois 表示，
可以由模形式的 Hecke 特征值恢复。
这正是 Langlands 纲领的局部--整体哲学在最经典情形下的成功实现。
从这个角度说，
FLT 的证明像是一束强光，
把更大理论的轮廓短暂照亮了一次。

\textbf{引入之后能得到什么？}
对读者而言，
终章最宝贵的收获并不只是“知道 FLT 如何被证明”，
而是知道接下来该去哪里。
你可以沿着模性走向更一般的自守形式与 Langlands；
可以沿着 BSD 走向 Heegner 点、Gross--Zagier 与 Kolyvagin；
也可以沿着形变理论走向 Serre 猜想与 $p$-adic Langlands。
终章不是句号，
而是一个很高的观景台。

\begin{figure}[ht]
\centering
\begin{tikzpicture}[>=Stealth, font=\small, node distance=1.5cm]
  \node[draw, rounded corners, fill=blue!8, minimum width=3.3cm, minimum height=0.9cm] (mf) {模形式与自守表示};
  \node[draw, rounded corners, fill=green!10, right=3.2cm of mf, minimum width=3.3cm, minimum height=0.9cm] (ec) {椭圆曲线与代数几何};
  \node[draw, rounded corners, fill=orange!14, below=1.8cm of mf, minimum width=3.3cm, minimum height=0.9cm] (gr) {Galois 表示};
  \node[draw, rounded corners, fill=yellow!16, below=1.8cm of ec, minimum width=3.3cm, minimum height=0.9cm] (lf) {$L$-函数与特殊值};
  \node[draw, rounded corners, fill=purple!10, below=1.8cm of $(gr)!0.5!(lf)$, minimum width=3.6cm, minimum height=0.9cm] (lg) {Langlands 与 $p$-adic Langlands};
  \node[draw, rounded corners, fill=red!8, below=1.7cm of lg, minimum width=4.0cm, minimum height=0.9cm] (bsd) {BSD、Serre 猜想与更远的方向};

  \draw[<->, very thick] (mf) -- node[above] {模性} (ec);
  \draw[<->, very thick] (mf) -- node[left] {Deligne 表示} (gr);
  \draw[<->, very thick] (ec) -- node[right] {Tate 模} (lf);
  \draw[<->, very thick] (gr) -- node[below] {局部--整体} (lf);
  \draw[->, very thick] (mf) -- (lg);
  \draw[->, very thick] (ec) -- (lg);
  \draw[->, very thick] (gr) -- (lg);
  \draw[->, very thick] (lf) -- (lg);
  \draw[->, very thick] (lg) -- (bsd);
\end{tikzpicture}
\caption{后记中的大图景：模性定理不是终点，而是把模形式、Galois 表示、代数几何与 $L$-函数连入 Langlands 网络的一个成功范例。}
\label{fig:ch10-langlands-network}
\end{figure}


\textbf{2001 年：完整模性定理。}
Wiles 与 Taylor--Wiles 处理的是半稳定椭圆曲线。
Breuil、Conrad、Diamond、Taylor 在 2001 年完成了一般情形，
从而证明：
\emph{每一条定义在 $\Q$ 上的椭圆曲线都是模的。}
这使得谷山--志村猜想真正成为“模性定理”。
从此以后，
第\ref{ch:eichler-shimura}章出现的解析对象与第\ref{ch:forms-curves}章开始就陪伴我们的椭圆曲线，
在有理数域上再也不是偶然相遇的两类对象，
而是同一算术现实的两种坐标。

\textbf{BCDT 证明思路。}
CSS 文集的历史终点，在今天看来正好接到 BCDT 的完整证明上。
其关键不是推翻 Wiles 的方法，
而是把局部形变条件大幅推广：
对非半稳定素数，需要更精细地分析有限平坦群方案、局部 Weil--Deligne 型与可允许形变环；
再把这些更复杂的局部条件重新接入 Taylor--Wiles 比较。
换句话说，BCDT 做的不是另起炉灶，
而是把“$R=T$ + 局部形变理论”真正扩展到全部椭圆曲线。

\textbf{Langlands 纲领。}
如果把模性定理放大看，
它说的是：
二维奇 Galois 表示与 $\GL_2/\Q$ 上的自守形式之间存在深刻对应。
Langlands 纲领则把这件事推广到一般 reductive 群与一般数域。
在那个更大的世界里，
模形式只是自守表示最熟悉的一角，
而谷山--志村定理是最早、也最鲜明的成功案例之一。
它告诉我们，
“把几何对象的 $L$-函数解释成自守 $L$-函数”并非幻想，
而是已经发生过的事实。

\textbf{BSD 猜想的现状。}
第\ref{ch:eichler-shimura}章提到 BSD 猜想时，
我们把它视为一个关于椭圆曲线有理点与 $L(E,s)$ 中心零点阶数的惊人猜想。
模性定理使得这个猜想的解析一侧能够真正接入自守方法。
Gross--Zagier 的公式与 Kolyvagin 的 Euler 系统因此得到应用，
从而在解析秩 $\le 1$ 的情形下，BSD 的许多关键结论已经成为定理。
也就是说，Wiles 之后，
模性不只拿下了 FLT，
还成为研究椭圆曲线秩与有理点的基础设施。

\textbf{Fontaine--Mazur 猜想。}
在 Wiles 之后，人们立刻会问一个更普遍的问题：
哪些 $p$-adic Galois 表示来自几何，或者更准确地说，来自自守对象？
Fontaine--Mazur 猜想给出的答案是：
若一个 $p$-adic 表示几乎处处无分歧，且在 $p$ 处是 de Rham（更强时是几何）并满足奇性条件，
那么它应来自模形式或更一般的自守表示。
从这个角度看，Wiles 定理只是这个巨大猜想在二维、奇、数域 $\Q$ 情形中的最经典样板。

\textbf{Serre 猜想。}
Serre 在 1987 年提出：
每个奇、不可约的二维 mod $p$ Galois 表示都应来自某个模形式。
这看上去像是 Wiles 定理的 residual 版本，
但其实方向更宽。
Khare 与 Wintenberger 在 2009 年证明了这一猜想，
这意味着 Frey 曲线当年触发的那条直觉——
“反常 residual 表示应该被模形式空间约束”——
最终在最一般的二维情形下也成立了。
它也反过来说明：
Wiles--Taylor--Wiles 技术并不是一次性工具，
而是后来一整代模性提升定理的模板。

\textbf{$p$-adic Langlands。}
Wiles 证明把 $p$-adic 表示推到舞台中央之后，
更进一步的问题自然出现：
能否直接把 $p$-adic Galois 表示与 $p$-adic 自守对象对应起来？
Breuil、Colmez、Emerton 等人的工作，
发展出了 $\GL_2(\Q_p)$ 的 $p$-adic Langlands 纲领。
在这个对应里，Banach 表示、$(\varphi,\Gamma)$-模与 completed cohomology 成为新的主角；
它们把局部 $p$-adic Hodge 理论与自守表示论直接粘合起来。
这条道路仍在快速发展，
其气质与终章非常相像：
一边是局部表示论，
另一边是几何与算术，
中间靠极其精细的同调与局部环技术搭桥。

\textbf{给继续学习的读者。}
若想沿着本章继续前进，
我会按下面的路线读书：
先回到 Cornell--Silverman--Stevens 的文集，
直接读 Wiles、Taylor--Wiles、Ribet 的原始叙述，
把终章中的“黑箱”逐个打开；
再读 Diamond--Shurman 的后续材料，
巩固模曲线与新形式理论；
再向数域与 Galois 方向走，
读 Neukirch 的 \emph{Algebraic Number Theory} 与 Serre 的 \emph{Abelian $\ell$-adic Representations and Elliptic Curves}；
若对 Langlands 好奇，
可以从 Gelbart、Bump 或 Buzzard--Gee 等综述材料开始。

写到最后，我很想把终章的精神浓缩成一句话：
\emph{费马大定理之所以伟大，
并不只是因为它终于被证明，
更因为它迫使数学家建立了一整套比原问题大得多的理论。}
从 CSS 到 BCDT，再到 Khare--Wintenberger 与 $p$-adic Langlands，
这张地图一直在继续扩展：
二维表示的模性、局部--整体对应、Selmer 群与形变环，
都不再是证明 FLT 的附属品，
而是现代算术几何的基本语言。
若本书能让读者在合上最后一页时感到——
“我不仅知道结论，而且看见了那张把模形式、椭圆曲线、Galois 表示与 $L$-函数连成一体的地图”，
那么这趟旅程就算值得。
